
There are many ways to create subgroups of a group.
Table~\ref{tab:subgroups} lists some commands to do so.
\orbitertableofcommands{subgroups}






\bigskip



We start with an example of an explicit permutation 
group using a known base and strong generating set, 
using the \verb'bsgs' command.
Here is the cyclic group of order 13 acting on the permutation domain $[0,12]$. 
The base is $(0)$. When creating a group, we supply a label in ascii text and in tex. 
Then we specify the degree of the action, and the group order. 
After that, we specify the number of generators and the generators themselves.
The labels will be used in reports about the group, for instance. 

\sourcecodevariable{GEN_C13}


\sourcecode{C13}

The makefile variable \verb'GEN_C13' is used to define the generator of the group, which is 
the cycle 
$$
(0,1,2,3,4,5,6,7,8,9,10,11,12).
$$
The command creates the group from the known base $0$. 
After that, several activities are invoked.
Specifically, these are group theoretic activities. 
They will be discussed in more detail in Section~\ref{sec:group:theoretic:activities}.

\bigskip

Let us take a closer look at the three activities performed in this example.
The \verb'-export_orbiter' command 
exports the group in Orbiter makefile format.
The file \verb'C13.makefile' is generated,
which can be used to recreate the permutation group in an Orbiter makefile. 
Here is the content of the file:

\sourcecode{C13_generated}


\bigskip

The activity \verb'-report'
produces a report for the cyclic group, shown below:

\orbiteroutput{GRAPHICS/LATEX/report_subgroup_C13}

\bigskip


The command  \verb'-save_elements_csv'
creates a csv file containing all group elements. 
Each group element is listed one-by-one, using the list notation of permutations.
The csv file \verb'C13_elts.csv' has the following content:
\orbiteroutput{GRAPHICS/LATEX/C13_elts}



\bigskip



It is possible to create a permutation group as a subgroup of the symmetric group, 
using the known base for the symmetric group.
Because the base of the symmetric group is large, 
this way of creating the group 
is less efficient than creating the group with a known (small) base.
Here is an example. We create $C_{13}$ as a subgroup of $\Sym(13)$.

\sourcecode{C13_as_subgroup}


\bigskip




For instance, the command

\sourcecode{J1}

creates the first Janko group as a subgroup of $\PGL(7,11).$ 

\bigskip

The command 

\sourcecode{PGL_3_11_singer}

creates a subgroup of the Singer cycle of order 7.
The Singer cycle in $\GL(d,q)$ 
is a generator for a subgroup of order $q^d-1$. 
It induces an element of order $\frac{q^d-1}{q-1}$ 
on the associated projective geometry $\PG(d-1,q).$
The additional integer parameter $k$ after the $\verb'-singer'$ command  is used to 
create the subgroup of index $k$ of the Singer cycle.

\bigskip


The command 

\sourcecode{PGL_3_11_singer_and_frobenius}

creates a subgroup of index 19 of the Singer cycle of $\PG(2,11)$, 
extended by a group of order $3$ 
that arises from the field extension $\bbF_{11}^3$ over $\bbF_{11}.$ 
The group created by this command has order $21$.


\bigskip


The group $\Sym(3)$ can be realized 
as a group of mappings of $\bbR^*= \bbR \cup \{\infty\}.$ 
The six mappings are 
$$
z, \;
\frac{1}{1-z}, \;  
\frac{1-z}{z}, \;  
\frac{1}{z}, \;  
1 - z, \;
\frac{z}{1-z}.   
$$
We can find a projective representation acting on $\PG(1,3).$ 
The elements of $\PG(1,3)$ have the form 
$$
\left\{
\left[
\begin{array}{c}
z\\
1\\
\end{array}
\right], 
z \in \bbF_3 \right\} \cup 
\left\{
\left[
\begin{array}{c}
1\\
0\\
\end{array}
\right]
\right\}. 
$$
The representing elements in $\PGL(2,3)$ are 
$$
\left[
\begin{array}{cc}
1&0\\
0&1\\
\end{array}
\right], \;
\left[
\begin{array}{cc}
0&1\\
2&1\\
\end{array}
\right], \;
\left[
\begin{array}{cc}
1&2\\
1&0\\
\end{array}
\right], \;
\left[
\begin{array}{cc}
0&1\\
1&0\\
\end{array}
\right], \;
\left[
\begin{array}{cc}
2&1\\
0&1\\
\end{array}
\right], \;
\left[
\begin{array}{cc}
1&0\\
1&2\\
\end{array}
\right].
$$
Using the second and fourth matrix as generators, encoded as 
$$
0,1,2,1, \quad 0,1,1,0,
$$
the following Orbiter command creates the group and produces the group table:

\sourcecode{Sym3_in_PGL_2_3_group_table}

The group table is:

\orbiteroutput{GRAPHICS/WRAPPERS/Sym3_in_PGL_2_3}

\bigskip

The next command creates a report of the group and the tree of elements, 
as well as a list of the group elements: 


\sourcecode{Sym3_in_PGL_2_3_report}


The tree of elements is:

\orbiteroutput{GRAPHICS/WRAPPERS/Sym3_in_PGL_2_3_tree}

The list of elements is:

\orbiteroutput{GRAPHICS/LATEX/Sym3_in_PGL_2_3_elements}


\bigskip


The quaternion group is a group of order $8$ generated by 
the following matrices over $\bbR$:
$$
i = 
\left[
\begin{array}{rr}
1&1\\
1&-1\\
\end{array}
\right], \quad
j = 
\left[
\begin{array}{rr}
-1&1\\
1&1\\
\end{array}
\right], \quad
k = 
\left[
\begin{array}{rr}
0&-1\\
1&0\\
\end{array}
\right].
$$
It is isomorphic to a subgroup of $\SL(2,3).$
The Orbiter command

\sourcecode{quaternion}

creates the group. 
The command produces the list of group elements shown below.

\orbiteroutput{GRAPHICS/LATEX/report_subgroup_quaternion}

The group table is created as csv file:

\orbiteroutput{GRAPHICS/LATEX/report_group_quaternion_table}


\bigskip

The group of the cube can be created over the field $\bbF_3$:

\sourcecode{cube_group}

The tetrahedral subgroup can be created as well:

\sourcecode{tetra_group}

The Hesse group of order 216 extended by the automorphism group of the field 
can be created in $\PG(3,4)$

\sourcecodevariable{GENERATORS_HESSE_GROUP}



\sourcecode{Hesse_group}

The group has order 432.

\bigskip



The Weyl group of type $E_8$ can be generated as a subgroup of $\GL(8,3)$ 
using the following command:

\sourcecodevariable{GENERATORS_WEYL_GROUP_E8}


\sourcecode{Weyl_E8}

A latex report is generated in the file \verb'GL_8_3_Subgroup_Weyl_E8_696729600_report.tex'.
This command uses generators found by Gabi Nebe:
$$
\verb'http://www.math.rwth-aachen.de/~Gabriele.Nebe/LATTICES/E8.b.html'
$$


\bigskip


We can test if a group is a subgroup of another. 
In the following example, we test whether 
$\PGO^+(6,2)$ is a subgroup of $\PSp(6,2)$. 
The fact that it is depends on the choice of forms 
associated with the groups and on the fact that the characteristic is two.

\sourcecode{test_subgroup}

Since the subgroup index is small (36), we create a set of 
coset representatives using the following command:

\sourcecode{coset_reps}

The coset representatives are written to a csv file.
The (shortened) list of coset representatives in latex is:
\orbiteroutput{GRAPHICS/LATEX/report_group_coset_reps_36}

\bigskip


The following command reads the vector 
of coset representatives from the file just created.

\sourcecode{coset_reps_read}




\bigskip


The following command creates the group 
$\PSp(6,2)$ and computes a point stabilizer subgroup.

\sourcecode{SP_6_2_point_stab_subgroup}


\bigskip

The next command creates the group 
$\PGO^+(6,2)$ and produces a report.

\sourcecode{PGOp_6_2_report}


\bigskip



The next command creates the group 
$\PGO^-(6,2)$ and produces a report.

\sourcecode{PGOm_6_2_report}


\bigskip

The next command creates the group 
$\PGO^-(6,2)$ and exports the group to Magma.

\sourcecode{PGOm_6_2_export_magma}


\bigskip



The next command creates the group 
$\PGO^+(6,2)$ and computes a point stabilizer subgroup.

\sourcecode{PGOp_6_2_point_stab_subgroup}


\bigskip

We can save the generators of $\PGO^+(6,2)$ 
from the report to a makefile variable:

\sourcecodevariablesmall{PGOp_6_2_GENS}

Using this variable, we can use the following command 
to recreate the group as a subgroup of 
$\PGL(6,2)$ and write a report:

\sourcecode{PGOp_6_2_linear}

In the following command, 
we compute a point stabilizer subgroup:

\sourcecode{PGOp_6_2_linear_stab_6}


\bigskip



